\chapter{Stochastic Processes}
\section{Markov Chains}
\noindent\textbf{Key Ideas and Methods.}\enspace
A Markov state must retain the history that affects the next transition: useful suffixes for pattern searches, task types for recursive work, and machine status for repairs. In continuous time, include the holding time as well as the next destination. Write first-step equations with their absorbing boundaries. A uniform positive chance of completion within a fixed number of steps gives a geometric tail and justifies moment calculations.

A stationary law must satisfy balance and normalization. For a finite irreducible chain it determines long-run occupation fractions; convergence of the distribution requires additional conditions, such as aperiodicity. Clearing or overwriting mechanisms can instead give direct couplings. For counts over a fixed window, retain covariances; for coverage, consider inclusion--exclusion and tail sums.
\par\smallskip



\problem{R01}{An Alarm in a Moving Window}{State reduction}
Independent trials succeed with probability $p\in(0,1)$ and fail with probability $q=1-p$. Starting with no observations, an alarm sounds as soon as the last at most three trials contain two successes. Thus either HH or HTH triggers it. Find the mean alarm time and the probability that the final trigger is HH.
\solution Before the alarm, retain three states: 0 means no H in the last two trials, A means the last trial is H, and B means the suffix is HT. The empty record also has state 0. Write the remaining means as $e_0,e_A,e_B$. First-step analysis gives
\[
 e_0=1+qe_0+pe_A,\quad e_A=1+qe_B,\quad e_B=1+qe_0.
\]
Two consecutive H's always finish the task, so the waiting time has a geometric tail and these means are finite. Substitution yields
\[
 (1-q^2)e_0=\frac1p+1+q,\qquad
 \boxed{e_0=\frac{1/p+2-p}{p(2-p)}}.
\]
Let $u_0,u_A,u_B$ denote the probabilities of an HH trigger. Then $u_0=u_A$, $u_A=p+qu_B$, and $u_B=qu_0$, giving
\[
 \boxed{u_0=\frac{p}{1-q^2}=\frac1{2-p}}.
\]
For fair trials the mean is $14/3$ and the HH probability is $2/3$.
\insight One failure does not erase all useful information. Recording only the current success-run length would confuse HT with an empty useful prefix.

\problem{P16}{Collecting Both Repeated Patterns}{Pattern coverage}
Toss a fair coin independently until the record has contained both HH and TT. Patterns may overlap, and the record is never cleared. Find the mean and variance of the total number $T$ of tosses.
\solution First wait for the first equal adjacent pair. After the initial toss, indicators of matching the preceding toss are independent fair Bernoulli variables: the first toss and the switch indicators uniquely determine the entire coin string, and every string is equally likely. Thus this stage has length $D=1+G$, with $G$ geometric of parameter $1/2$ on the positive integers. Therefore $\E D=3$ and $\Var D=2$.

After HH first appears, wait for TT; after TT first appears, the symmetric task remains. Let $U$ be this second-stage length. States 0 and 1 mean respectively no trailing target symbol and one trailing target symbol. Their means obey
\[
 a_0=1+\tfrac12a_0+\tfrac12a_1,\qquad
 a_1=1+\tfrac12a_0,
\]
so $a_0=6,a_1=4$. Two consecutive target symbols give a geometric tail, ensuring all moments are finite. Second moments satisfy
\[
 b_0=1+a_0+a_1+\tfrac12b_0+\tfrac12b_1,
 \quad b_1=1+a_0+\tfrac12b_0.
\]
Substitution gives $b_0=58$, hence $\Var U=58-36=22$. Future tosses are independent of the first-stage history. By H--T symmetry, the conditional second-stage law depends neither on the first pattern nor on $D$, proving independence of $D,U$. Consequently,
\[
 \boxed{\E T=9},\qquad \boxed{\Var T=24}.
\]
\insight Collecting both patterns is a coverage condition. Their two waiting times from an empty record cannot simply be added; the first completion determines the suffix at the start of the second stage.


\problem{R02}{A Queue with Random Clearing}{Resetting and coupling}
A queue length $X_n\in\{0,1,\ldots\}$ independently receives one of three operations at each step: increase by 1 with probability $a$, decrease by 1 with probability $b$ (remain at zero if empty), or clear the queue with probability $c$. Assume $a,b,c>0$ and $a+b+c=1$. Find the stationary law and mean length, and bound convergence to stationarity in total variation from any initial state.
\solution At $k\ge1$, stationarity requires $\pi_k=a\pi_{k-1}+b\pi_{k+1}$. Put
\[
 r=\frac{1-\sqrt{1-4ab}}{2b}\in(0,1),
 \qquad \boxed{\pi_k=(1-r)r^k}.
\]
Indeed, $br^2-r+a=0$. This polynomial is positive at 0 and equals $-c$ at 1, so the chosen root lies in $(0,1)$. Substitution verifies the positive-state equations. The zero-state equation is $\pi_0=b(\pi_0+\pi_1)+c$, which follows from the same root identity. The geometric series normalizes the probabilities. Consequently,
\[
 \boxed{\E_\pi X=\frac r{1-r}}.
\]
Couple the chain from an arbitrary state with a stationary chain by using the same operations. After the first clearing, they coincide forever. The probability of no clearing in the first $n$ steps is $(1-c)^n$, so the coupling inequality gives
\[
 \boxed{\|\mathcal L(X_n)-\pi\|_{\rm TV}\le(1-c)^n}.
\]
This also proves uniqueness. The chain is irreducible and positive recurrent, so the mean positive-time return to 0 is $1/\pi_0=1/(1-r)$; clearing or attempting service at an empty queue counts as a return.
\insight Clearing erases the initial condition regardless of queue size. The stationary tail parameter and the convergence bound describe different properties.

\problem{R03}{Work Generated by Two Task Types}{Multitype recursion}
After processing an A task, with equal probabilities it produces either no new tasks or one A and one B. A B task produces no new tasks with probability $2/3$ and two A tasks with probability $1/3$. All choices are independent. Processing one task counts as one operation. Starting from one A, prove that all tasks are eventually processed almost surely, and find the expected total number of each type. If their unit costs are $c_A,c_B\ge0$, find the expected total cost.
\solution Assign weights 9 and 7 to pending A and B tasks. Processing A produces expected child weight $(9+7)/2=8$, while processing B produces $2\cdot9/3=6$. Both reduce expected pending weight by 1. Let $T$ be the completion time. Nonnegativity of the weight after $n\wedge T$ operations and successive conditional expectations give
\[
 \E(n\wedge T)\le9.
\]
Monotone convergence yields $\E T\le9$, hence $T<\infty$ almost surely. The total counts are therefore integrable. Denote their means by $A,B$. Counting all children of each type gives
\[
 A=1+\tfrac12A+\tfrac23B,\qquad B=\tfrac12A.
\]
Thus $A=1+5A/6$, and
\[
 \boxed{A=6,\quad B=3,\quad \E T=9},\qquad
 \boxed{\E C=6c_A+3c_B}.
\]
\insight Offspring type affects future work. A single unweighted offspring average loses this information; positive weights establish integrability before the count equations are solved.

\problem{R04}{Competing Consensus States}{Birth--death chain}
There are $m\ge2$ binary bits. Each step chooses one uniformly and flips it. Initially $k$ bits are 1, where $1\le k<m$. Stop when all bits are 0 or all are 1. Find the probability $h_k$ of reaching all 1's first. Also find the mean stopping time for $m=4,k=2$.
\solution The number of 1's is sufficient. From any nonabsorbing state, a prescribed sequence that flips each existing 1 to 0 reaches all 0's within $m$ steps with probability at least $m^{-m}$. Hence the stopping time has a geometric tail. For $1\le k<m$,
\[
 h_k=\frac{k}{m}h_{k-1}+\frac{m-k}{m}h_{k+1},
 \qquad h_0=0,\quad h_m=1.
\]
Let $d_k=h_k-h_{k-1}$. Then $(m-k)d_{k+1}=kd_k$, so
\[
 d_k=\frac{d_1}{\binom{m-1}{k-1}},\qquad
 \boxed{h_k=\frac{\displaystyle\sum_{j=0}^{k-1}\binom{m-1}{j}^{-1}}
 {\displaystyle\sum_{j=0}^{m-1}\binom{m-1}{j}^{-1}}}.
\]
For $m=4$, symmetry gives remaining means $e_1=e_3$, and
\[
 e_1=1+\tfrac34e_2,\qquad e_2=1+e_1.
\]
Thus $e_1=7$ and $\boxed{e_2=8}$. The absorption probabilities are $h_1=3/8,h_2=1/2,h_3=5/8$.
\insight The interior drift points toward the middle, so absorption probabilities are generally not $k/m$. The harmonic function must reflect the actual transitions.

\problem{R05}{A Race Between Overlapping Patterns}{Prefix automaton}
Independent signals are H with probability $p\in(0,1)$ and T with probability $q=1-p$. Starting with no observations, stop on the first occurrence of substring HHT or HTH. Find the probability that HHT appears first and the mean stopping time for fair signals.
\solution Keep the longest suffix that is a prefix of either target. The nonterminal states are $0,H,HH,HT$. In particular, H from HH leaves state HH, while T from HT returns to 0. Let $u_s$ be the probability that HHT wins. Then
\[
 u_0=u_H,\quad u_{HH}=1,\quad u_{HT}=qu_0,
 \quad u_H=p+qu_{HT}.
\]
From HH a T eventually occurs because $q>0$, explaining $u_{HH}=1$. Solving gives
\[
 \boxed{u_0=\frac1{2-p}}.
\]
For fair signals, write $e_s$ for the remaining means. Independent blocks of length 3 equal HHT with probability $1/8$, ensuring finiteness. First-step equations are
\[
 e_0=2+e_H,\quad e_{HH}=2,\quad
 e_{HT}=1+\tfrac12e_0,\quad
 e_H=1+\tfrac12e_{HH}+\tfrac12e_{HT}.
\]
Hence $e_0=9/2+e_0/4$, giving $\boxed{e_0=6}$.
\insight Prefix overlaps determine what information survives a failed match. Equal pattern lengths do not imply equal winning probabilities.

\problem{R06}{Counting Runs in Markov Signals}{A fixed window and covariance}
A stationary two-state chain takes values H and T, with $P(H,T)=a$ and $P(T,H)=b$, where $0<a,b<1$. Observe $n\ge2$ consecutive signals. A run is a maximal consecutive string of equal signals. Find the mean and variance of the number $R_n$ of runs, and simplify when $a=b$.
\solution Let $I_j=\ind_{\{X_{j+1}\ne X_j\}}$ and $m=n-1$. Then $R_n=1+\sum_{j=1}^m I_j$. The stationary law is $(b/(a+b),a/(a+b))$. Put
\[
 r=\frac{2ab}{a+b},\quad \lambda=1-a-b,\quad
 C=\frac{ab(a-b)^2}{(a+b)^2}.
\]
Thus $\E I_j=r$. Conditional on a switch, the new state is equally likely to be H or T, since both stationary directional flows equal $ab/(a+b)$. After $k-1$ further steps, the deviation in the next-switch probability is multiplied by $\lambda^{k-1}$. Therefore
\[
 \Prob(I_{j+k}=1\mid I_j=1)
 =r+\lambda^{k-1}\left(\frac{a+b}{2}-r\right),
\]
so $\Cov(I_j,I_{j+k})=C\lambda^{k-1}$. Summing gives
\[
 \boxed{\E R_n=1+mr},
\]
\[
 \boxed{\Var R_n=mr(1-r)
 +2C\frac{m-1-m\lambda+\lambda^m}{(1-\lambda)^2}}.
\]
The fraction is the sum $\sum_{k=1}^{m-1}(m-k)\lambda^{k-1}$ and equals zero when $m=1$. If $a=b$, independent Bernoulli variables with parameter $a$ can determine the switches, so $R_n-1\sim\operatorname{Bin}(n-1,a)$. This agrees with $C=0$.
\insight Counts over a fixed window require covariances. Knowing the one-step switch rate alone does not establish a binomial law.

\problem{P15}{Exact Mixing in a Shift Register}{Finite memory}
A register retains $r\ge1$ binary bits, including leading zeroes. Each step discards the leftmost bit, shifts the rest left, and appends an independent bit that equals 1 with probability $p\in(0,1)$. As an integer, the state obeys $S_{n+1}=(2S_n+\xi_{n+1})\bmod 2^r$. Find its stationary law. Starting from all zeroes, find the total variation distance to stationarity after $n$ steps.
\solution After $r$ updates, all old bits have been overwritten, so the state consists of the last $r$ independent bits. If $w(s)$ counts the 1's in state $s$, then
\[
 \boxed{\pi(s)=p^{w(s)}(1-p)^{r-w(s)}}.
\]
Every initial state has exactly this law after $r$ steps, proving stationarity and uniqueness. For $n<r$, the all-zero start leaves the first $r-n$ bits fixed at zero, while the last $n$ bits are independent fresh bits. The law is therefore $\pi$ conditioned on $A_n=$ ``the first $r-n$ bits are zero.'' For an event $A$ with $\pi(A)>0$, the total variation distance between $\pi$ and its conditional law on $A$ is $1-\pi(A)$: sum the absolute probability differences separately on $A$ and its complement. Hence
\[
 \boxed{d_{\rm TV}(\mathcal L(S_n),\pi)=
 \begin{cases}1-(1-p)^{r-n},&0\le n<r,\\0,&n\ge r.\end{cases}}
\]
\insight The initial information is completely overwritten, yielding exact mixing in finitely many steps. The stationary law is uniform only when $p=1/2$.


\problem{R07}{Two Machines and One Repairer}{Continuous-time chain}
Each of two working machines fails independently at rate $\lambda>0$. A single repairer works whenever any machine is broken. Each repair takes an independent exponential time of rate $\mu>0$; the repaired machine immediately returns to service. All failure and repair times are independent. The system is available if at least one machine works. Find its long-run availability and the mean time from both working to the first time both are broken.
\solution Let $K_t=0,1,2$ count broken machines. The upward transition rate is $(2-k)\lambda$, and the downward rate is $\mu$. Adjacent balance gives
\[
 2\lambda\pi_0=\mu\pi_1,\qquad
 \lambda\pi_1=\mu\pi_2.
\]
With $D=\mu^2+2\lambda\mu+2\lambda^2$, normalization yields
\[
 (\pi_0,\pi_1,\pi_2)=\frac1D(\mu^2,2\lambda\mu,2\lambda^2),
 \qquad \boxed{\text{availability}=1-\frac{2\lambda^2}{D}}.
\]
The finite irreducible continuous-time chain has these long-run proportions. Let $e_0,e_1$ be the remaining means before state 2. Exponential holding times and the first jump give
\[
 e_0=\frac1{2\lambda}+e_1,\qquad
 e_1=\frac1{\lambda+\mu}+\frac\mu{\lambda+\mu}e_0.
\]
Rearranging gives
\[
 \boxed{e_0=\frac{3\lambda+\mu}{2\lambda^2}}.
\]
Finiteness follows from a uniform positive probability, over a fixed time block, of enough failures to reach state 2 without a completed repair.
\insight With one repairer, the total repair rate remains $\mu$ even when both machines are broken. The machines cannot be treated as independent two-state systems.

\problem{R08}{Hitting Times on a Weighted Triangle}{First-step equations}
A triangle with vertices A, B, C has positive edge weights $w_{AB}=a,w_{BC}=b,w_{CA}=c$. At each step, choose a neighbor in proportion to the incident edge weights. Find the mean hitting times from A to B and from B to A, and their mean round-trip duration. With $a=b=1$, does increasing $c$ speed up the trip from A to B?
\solution Let $u_A,u_C$ be mean times to B. First-step analysis gives
\[
 u_A=1+\frac c{a+c}u_C,\qquad
 u_C=1+\frac c{b+c}u_A.
\]
The finite irreducible chain has finite hitting means. Put $D=ab+ac+bc$. Eliminating $u_C$ yields
\[
 \boxed{\E_A T_B=\frac{(a+2c)(b+c)}D}.
\]
Interchanging A and B interchanges $b,c$. If $C$ denotes the time to reach B and subsequently return to A, the strong Markov property gives
\[
 \boxed{\E_B T_A=\frac{(a+2b)(b+c)}D},\qquad
 \boxed{\E_A C=\frac{2(a+b+c)(b+c)}D}.
\]
When $a=b=1$, the first mean simplifies to $\boxed{1+c}$. Increasing the A--C weight makes reaching B slower, because the walk spends more steps shuttling between A and C.
\insight Increasing one edge weight also decreases the probabilities of other exits. Graph distance alone does not determine the duration of a weighted random walk.

\problem{R09}{Covering a Reflecting Segment}{Decomposition at a boundary}
A walk lives on $\{0,1,\ldots,N\}$, where $N\ge2$. Interior steps move left or right with equal probability. From 0 the next step must be to 1, and from $N$ it must be to $N-1$. Starting from $1\le i<N$, stop when both endpoints have been visited; call the time $C$. Find the probability that $N$ is visited last and find $\E_iC$.
\solution Let $\sigma$ be the first visit to either endpoint. Before this time the walk is the usual absorbing interval walk. Its harmonic and mean-time equations give
\[
 \Prob_i(X_\sigma=0)=\frac{N-i}{N},\qquad
 \E_i\sigma=i(N-i).
\]
For example, $v_i=i(N-i)$ satisfies $v_i=1+(v_{i-1}+v_{i+1})/2$ and zero endpoint values. Visiting $N$ last is exactly the event of visiting 0 first.

Let $u_j$ be the mean time from $j$ to $N$, retaining reflection at 0. Then $u_N=0$, $u_0=1+u_1$, and the interior mean-time equation is unchanged. The quadratic $u_j=N^2-j^2$ satisfies all these conditions. Hence a full crossing from 0 to $N$ has mean $N^2$, as does the reverse crossing. The strong Markov property now gives
\[
 \boxed{\Prob_i(\text{last endpoint }N)=\frac{N-i}{N}},\qquad
 \boxed{\E_i C=i(N-i)+N^2}.
\]
A uniform positive probability of crossing within a fixed number of steps ensures integrability of these times. Visiting both endpoints necessarily visits every intermediate vertex.

\problem{R10}{Coverage by Pair Sampling}{Inclusion--exclusion and tail sums}
There are $n\ge3$ labels. Each round independently chooses a uniformly random unordered pair of distinct labels. Stop when all labels have appeared, and let $T$ be the number of rounds. Give finite sums for its tail probability, mean and second moment, then compute the mean and variance for $n=4$.
\solution For a specified set of $k$ labels, the probability that a round avoids it is
\[
 q_k=\frac{\binom{n-k}{2}}{\binom n2},
\]
where the numerator is zero if its upper argument is less than 2. Failure to cover after $t$ rounds means at least one label is absent. Inclusion--exclusion gives
\[
 \Prob(T>t)=\sum_{k=1}^n(-1)^{k+1}\binom nk q_k^t,
 \qquad t=0,1,\ldots,
\]
with $0^0=1$ at $t=0$. Every $q_k<1$, so each geometric series in this finite sum converges. Integer tail-sum identities yield
\[
 \boxed{\E T=\sum_{k=1}^n(-1)^{k+1}\binom nk\frac1{1-q_k}},
\]
\[
 \boxed{\E T^2=\sum_{k=1}^n(-1)^{k+1}\binom nk
 \frac{1+q_k}{(1-q_k)^2}}.
\]
For $n=4$, $(q_1,q_2,q_3,q_4)=(1/2,1/6,0,0)$. Substitution gives
\[
 \boxed{\E T=\frac{19}{5}},\quad
 \E T^2=\frac{423}{25},\quad
 \boxed{\Var T=\frac{62}{25}}.
\]
\insight One round may discover two labels, so successive discovery stages are not separate geometric waits. Events specifying missing labels are better suited to inclusion--exclusion.

\par\Needspace{8\baselineskip}
\subsection*{Section Exercises}
\exercise{OR01}{Waiting for Consecutive Heads}
A fair coin is tossed independently until $n$ consecutive heads have appeared, where $n\ge1$ is an integer. Find the expected total number $T_n$ of tosses.

\exercise{OR07}{The Long-Run Frequency of Getting Wet}
A commuter owns $n\ge1$ umbrellas. Initially, $k$ are at home and $n-k$ are at the office, where $0\le k\le n$. Each morning the commuter travels from home to the office and returns in the evening. Independently on each trip, it rains with probability $p\in[0,1]$. An umbrella is carried only when it rains and one is available at the departure location. The commuter gets wet if it rains and no umbrella is available. Find the long-run fraction of trips on which the commuter gets wet.


\section{Random Walks}
\noindent\textbf{Key Ideas and Methods.}\enspace
Specify absorption, reflection and restarting before writing a walk's equations. If direction has memory or the boundary moves with past records, include direction or drawdown in the state. Hitting probabilities satisfy harmonic equations, whereas mean durations satisfy first-step equations with an added time cost. Conditioning on success reweights transitions by subsequent success probabilities; merely deleting the unsuccessful endpoint does not define the conditioned walk.

Martingales offer another route to boundary calculations, but optional stopping needs hypotheses. First stop at $\tau\wedge n$, then justify passage to the limit through uniform integrability or an integrable dominating variable. Almost-sure termination alone is insufficient. Jumps may also overshoot a target, so use the actual terminal value rather than automatically substituting the boundary level.
\par\smallskip



\problem{R11}{Emptying with Random Restarts}{Resetting and optimization}
An integer process starts at 1 and stops at 0. At each nonterminal step it resets to 1 with probability $\rho\in(0,1)$; otherwise it moves left or right with equal probability. A reset counts as a step, and all choices are independent. Find the mean stopping time and the choice of $\rho$ minimizing it.
\solution From any positive state, two successive operations ``reset, then move left'' finish with probability $\rho(1-\rho)/2$. Thus the remaining means have a uniform finite bound. Writing $e_i$ for the mean from $i$ gives
\[
 e_i=1+\rho e_1+\frac{1-\rho}{2}(e_{i-1}+e_{i+1}),
 \qquad e_0=0.
\]
The characteristic roots are reciprocal. The smaller one is
\[
 r=\frac{1-\sqrt{1-(1-\rho)^2}}{1-\rho}.
\]
A constant particular solution is $(1+\rho e_1)/\rho$. Boundedness discards the larger root, and the zero boundary gives
\[
 e_i=\frac{1+\rho e_1}{\rho}(1-r^i),\qquad
 \boxed{e_1=\frac{1-r}{\rho r}}.
\]
The root equation gives $1-\rho=2r/(1+r^2)$. Reparameterizing by $r\in(0,1)$,
\[
 \rho=\frac{(1-r)^2}{1+r^2},\qquad
 e_1=\frac{1+r^2}{r(1-r)}.
\]
The derivative has the sign of $r^2+2r-1$, and the mean diverges at both endpoints. The unique minimizer is therefore $r=\sqrt2-1$, giving
\[
 \boxed{\rho_*=1-\frac1{\sqrt2}},\qquad
 \boxed{\min\E T=2+2\sqrt2}.
\]
\insight Restarting retrieves paths that have wandered far away but also wastes progress near the target. The optimal rate balances these two effects.

\problem{R12}{Occupation Counts on Successful Paths}{Conditioning and the Green function}
A fair integer walk starts at $i\in\{1,\ldots,N-1\}$ and stops at 0 or $N$, at time $\tau$. Conditional on absorption at $N$, find the expected number of visits to each interior state $j$ before stopping, counting the initial time, and hence find the conditional mean duration.
\solution First find the unconditional occupation function
\[
 G(i,j)=\E_i\sum_{t=0}^{\tau-1}\ind_{\{X_t=j\}}.
\]
For fixed $j$, it is harmonic and linear away from $i=j$. At $j$ it obeys $G(j,j)=1+[G(j-1,j)+G(j+1,j)]/2$, and it vanishes at both endpoints. Matching the two linear pieces and their slope jump at $j$ gives
\[
 G(i,j)=\frac{2\min(i,j)[N-\max(i,j)]}{N}.
\]
The success probability is $h_i=i/N$. Following a visit to $j$, the future success probability is $h_j$. Conditional expectation term by term in the nonnegative visit count gives success-weighted mean $h_jG(i,j)$. Dividing by $h_i$ yields
\[
 \boxed{G^*(i,j)=\frac{2j\min(i,j)[N-\max(i,j)]}{Ni}}.
\]
Split the sum into $j\le i$ and $j>i$, and use the sum of squares to obtain
\[
 \boxed{\E_i(\tau\mid X_\tau=N)=\sum_{j=1}^{N-1}G^*(i,j)
 =\frac{N^2-i^2}{3}}.
\]
Alternatively, the success-weighted duration $v_i=i(N^2-i^2)/(3N)$ satisfies $v_i=h_i+(v_{i-1}+v_{i+1})/2$ and zero boundaries. The geometric exit tail on a finite interval ensures all these expectations are finite.
\insight Conditioning reweights visits to $j$ by $h_j/h_i$. Removing the unsuccessful endpoint without changing the transition probabilities would not produce the conditioned law.

\problem{R13}{A Walk with Directional Memory}{An enlarged state}
A walk on $\{0,\ldots,N\}$ stops at an endpoint, where $N\ge2$. At an interior state, it keeps its preceding direction with probability $\alpha\in(0,1)$ or reverses it with probability $1-\alpha$, then moves one unit in the chosen direction. The initial position is $1\le i<N$, and the preceding direction is known. Find the probabilities $h_i^+,h_i^-$ of reaching $N$ first when that direction is right or left.
\solution Position alone is not a Markov state; direction must be included. The interior equations are
\[
 h_i^+=\alpha h_{i+1}^++(1-\alpha)h_{i-1}^-,\qquad
 h_i^-=(1-\alpha)h_{i+1}^++\alpha h_{i-1}^-.
\]
The only reachable absorbing states are $h_N^+=1,h_0^-=0$. Try $h_i^-=Ai$ and $h_i^+=Ai+B$. Substitution gives $(1-\alpha)B=(2\alpha-1)A$, while the upper boundary gives $AN+B=1$. With
\[
 D=N(1-\alpha)+2\alpha-1=(N-2)(1-\alpha)+1>0,
\]
we obtain
\[
 \boxed{h_i^- =\frac{i(1-\alpha)}D,\qquad
 h_i^+=\frac{i(1-\alpha)+2\alpha-1}D}.
\]
These candidates satisfy every equation. Consecutive rightward steps have a uniform positive probability of absorption within $N$ steps, so the finite-state boundary problem has a unique probability solution. When $\alpha=1/2$, both expressions reduce to $i/N$. When $\alpha>1/2$, a rightward preceding direction increases the success probability.
\insight Averaging left and right steps discards directional persistence. The initial direction is essential information.

\problem{R14}{A Triangle Walk with a Fixed Endpoint}{Markov bridge}
A walk on triangle A, B, C chooses either other vertex with equal probability and starts at A. Conditional on being at A after $n\ge2$ steps, find the probability of being at B at an integer time $0\le k\le n$, and the expected total number of visits to A from time 0 through $n$. Evaluate these for $n=4,k=2$.
\solution Put $r=-1/2$. The recurrence $a_{t+1}=(1-a_t)/2,a_0=1$ gives the return probability and the probability of any specified different vertex:
\[
 a_t=\frac{1+2r^t}{3},\qquad b_t=\frac{1-r^t}{3}.
\]
The Markov property at time $k$ and transition symmetry give
\[
 \boxed{\Prob(X_k=B\mid X_n=A)=\frac{b_kb_{n-k}}{a_n}}.
\]
For $L=\sum_{k=0}^n\ind_{\{X_k=A\}}$,
\[
 \E(L\mid X_n=A)=\frac1{a_n}\sum_{k=0}^n a_ka_{n-k}.
\]
Expanding the product as $[1+2r^k+2r^{n-k}+4r^n]/9$ and summing the geometric series yields
\[
 \boxed{\E(L\mid X_n=A)=
 \frac{(n+1)(1+4r^n)+4(1-r^{n+1})/(1-r)}{3(1+2r^n)}}.
\]
For $n=4,k=2$, $a_4=3/8,b_2=1/4$, so the probability is $\boxed{1/6}$ and the visit mean is $\boxed{8/3}$. The conditioned chain is generally time-inhomogeneous: at time $k$, its transition from $u$ to $v$ is multiplied by $P^{n-k-1}(v,A)/P^{n-k}(u,A)$, used only on conditionally reachable states with positive denominator.

\problem{R15}{Alternating Additions to a Threshold}{Parity of a renewal count}
Two players take turns generating independent $U[0,1]$ increments. All increments are added to one common total, initially zero. The player whose increment first makes the total exceed $t\in(0,1]$ wins. Find the first player's winning probability, the fair threshold, and the mean and variance of the total number $N$ of turns.
\solution For $n\ge1$, no crossing means that the first $n$ increments sum to at most $t$. Since $t\le1$, this is a simplex on which the individual upper bounds do not bind. Its volume is
\[
 \Prob(N>n)=\frac{t^n}{n!};\qquad \Prob(N>0)=1.
\]
This tail ensures all moments are finite. The probability generating function follows from the tail-sum identity:
\[
 G(z)=\E z^N=1+(z-1)\sum_{n=0}^\infty z^n\Prob(N>n)
 =1+(z-1)e^{zt}.
\]
The first player wins exactly when $N$ is odd. Therefore
\[
 \boxed{\Prob(\text{first player wins})=\frac{1-G(-1)}2=e^{-t}},\qquad
 \boxed{t_{\rm fair}=\log2}.
\]
Differentiation gives $G'(1)=e^t,G''(1)=2te^t$, hence
\[
 \boxed{\E N=e^t},\qquad
 \boxed{\Var N=(2t+1)e^t-e^{2t}}.
\]
\insight Monotone accumulation turns the winner into a question about the parity of a stopping count. The mean number of turns alone does not determine first-player advantage.

\problem{R16}{Drawdown from the Running Maximum}{A moving boundary}
A fair integer walk $S_n$ starts at 0, and $M_n=\max_{0\le j\le n}S_j$. For an integer $d\ge1$, stop at the first time $\tau$ when $M_n-S_n=d$. Find $\E\tau,\Var\tau$ and the distribution of the terminal running maximum $M_\tau$.
\solution The gap $Y_n=M_n-S_n$ is a finite-state chain. At 0 it stays or moves to 1 with equal probability; at $1,\ldots,d-1$ it moves one unit in either direction; $d$ is absorbing. A block of $d$ downward walk steps ends the process, so all stopping moments are finite. With $m_i=\E_i\tau$, we have $m_0=1+(m_0+m_1)/2$, the interior equation $m_i=1+(m_{i-1}+m_{i+1})/2$, and $m_d=0$. The quadratic solution is
\[
 m_i=D-i(i+1),\qquad D=d(d+1),\qquad \boxed{\E\tau=D}.
\]
Let $v_i$ be the second moment. Expanding $(1+\tau')^2$ gives
\[
 v_0-v_1=4D-2,\qquad
 v_{i-1}-2v_i+v_{i+1}=2-4m_i,\qquad v_d=0.
\]
Successive summation of first differences yields
\[
 v_0=\sum_{j=0}^{d-1}(d-j)(4m_j-2)=\frac{5D^2-D}{3},
 \qquad \boxed{\Var\tau=\frac{D(2D-1)}3}.
\]
For the maximum, restart at each new high $h$. The probability of reaching $h+1$ before $h-d$ is $d/(d+1)$. Repeating this competition by the strong Markov property gives
\[
 \boxed{\Prob(M_\tau\ge k)=\left(\frac d{d+1}\right)^k,
 \quad k=0,1,\ldots}.
\]
\insight Absolute position and the running maximum form a two-dimensional state. Their difference turns the moving boundary into reflection with a self-loop.

\problem{R17}{The Limit of Fractional Betting Wealth}{A nonnegative martingale}
Starting with wealth $x>0$, bet a fixed fraction $f\in(0,1)$ on independent fair rounds. Wealth is multiplied by $1+f$ after a win and by $1-f$ after a loss. Write it as $Z_n$. Show that $\E Z_n=x$ for every $n$, yet $Z_n\to0$ almost surely. For a target $A>x$, bound the probability of ever reaching or exceeding $A$, and explain why terminal wealth cannot simply be set equal to $A$.
\solution The conditional mean multiplier is 1, so $Z_n$ is a nonnegative martingale and $\E Z_n=x$ at every finite time. However, the strong law gives
\[
 \frac1n\log\frac{Z_n}{x}\longrightarrow
 \tfrac12\log(1+f)+\tfrac12\log(1-f)
 =\tfrac12\log(1-f^2)<0.
\]
Thus wealth tends to zero almost surely. Its constant mean shows that the martingale family is not uniformly integrable.

Let $\tau=\inf\{n:Z_n\ge A\}$. Bounded optional stopping gives $\E Z_{n\wedge\tau}=x$. The stopped family is bounded by $(1+f)A$: before hitting it is below $A$, and a crossing starts below $A$. Its limit is $Z_\tau$ on hitting and zero on nonhitting. Dominated convergence yields
\[
 x=\E[Z_\tau\ind_{\{\tau<\infty\}}].
\]
Using $A\le Z_\tau<(1+f)A$ gives
\[
 \boxed{\frac{x}{(1+f)A}\le\Prob(\tau<\infty)\le\frac xA}.
\]
\insight Fair arithmetic returns do not ensure growth on typical paths. Multiplicative jumps overshoot the target, so the terminal value is not a fixed boundary value.

\problem{R18}{Costs Until the First Collision}{Stopping and selection bias}
Each independent draw has type $j\in\{1,\ldots,m\}$ with probability $p_j>0$, where the probabilities sum to 1, and cost $c_j\ge0$. Stop at the first repeated type. Let $N$ include that final draw. Find the expected total cost and the distribution of the final type. Evaluate the case of probabilities $2/3,1/3$ and costs 1 and 4.
\solution Let $e_k(p)$ be the elementary symmetric polynomial of degree $k$ in $p_1,\ldots,p_m$, with $e_0=1$. The first $k$ draws are distinct with probability $k!e_k(p)$, so
\[
 \boxed{\E N=\sum_{k=0}^m k!e_k(p)}.
\]
Here $2\le N\le m+1$. If draw $t$ costs $C_t$ and $\bar c=\sum_jp_jc_j$, the event $N\ge t$ depends only on earlier draws and is independent of $C_t$. A finite sum gives
\[
 \boxed{\E\sum_{t=1}^N C_t=\bar c\,\E N}.
\]
For the final type to be $j$ at time $k+1$, the first $k$ distinct types must include $j$. Ordering these types and then drawing $j$ once more has probability $k!p_j^2e_{k-1}(p_{-j})$, where $p_{-j}$ omits entry $j$. Thus
\[
 \boxed{\Prob(J_{\rm last}=j)=p_j^2\sum_{k=1}^m k!e_{k-1}(p_{-j})}.
\]
For two types with $p=2/3,q=1/3$, $\E N=2+2pq=22/9$ and $\bar c=2$, so the mean total cost is $\boxed{44/9}$. The final type is the first one with probability $p^2(1+2q)=20/27$. The final draw's mean cost is $\boxed{16/9}$, not its unconditional single-draw mean 2.
\insight Cumulative cost uses a continuation event measurable from the past. The terminal type is selected by the stopping rule and does not satisfy the same independence argument.

\par\Needspace{8\baselineskip}
\subsection*{Section Exercises}
\exercise{OR11}{Hitting Zero with Asymmetric Jumps}
A random walk starts at a positive integer $n$. Its independent increments equal $-1$ with probability $p\in(0,1)$ and $+2$ with probability $1-p$. Find the probability $u_n$ that the walk ever reaches 0.

\exercise{OR16}{The Second Moment of an Exit Time}
Let $S_0=0$ and let $S_n$ be a random walk with independent increments, each equal to $+1$ or $-1$ with probability $1/2$. For positive integers $a,b$, define
\[
 \tau=\inf\{n\ge0:S_n=-a\text{ or }S_n=b\}.
\]
Find $\Prob(S_\tau=b)$, $\E\tau$, and $\E\tau^2$.


\section{Optimal Stopping}
\noindent\textbf{Key Ideas and Methods.}\enspace
Optimal stopping compares the payoff available now with the value of continuing and then acting optimally. The state must include the remaining opportunities, any offer that can be recalled, and the posterior for an unknown parameter. Do not charge costs already paid a second time. Compare feasible actions at each observation before averaging over observations; equality may allow several optimal decisions.

Backward induction verifies finite-horizon policies. With unlimited opportunities, establish a finite value, justify the limit of the recursion, and check that the proposed policy attains it. Discounting can provide a uniform bound on truncation error. Replication and superhedging also require nodewise checks of self-financing, borrowing and short-selling restrictions, not just comparison of terminal payoffs.
\par\smallskip



\problem{R19}{A Digital Payoff under Funding Constraints}{Replication and superhedging}
A stock starts at 100. In each of two periods its price is multiplied by 2 or $1/2$, and every path has positive probability. Cash earns zero interest; arbitrary fractional stock positions can be traded without friction. The terminal payoff is $\ind_{\{S_2\ge100\}}$. First allow borrowing and short selling, and find the exact replication cost and holdings. Then require both stock holdings and cash to be nonnegative after each rebalance. Find the smallest initial wealth covering the payoff on every path.
\solution At the up node with stock price 200, both terminal payoffs are 1, replicated by cash 1. At the down node with price 50, successor prices are 100 and 25. Holdings $(\Delta,B)$ solve
\[
 100\Delta+B=1,\qquad25\Delta+B=0,
\]
so $\Delta=1/75,B=-1/3$, with node value $1/3$. Initial holdings solve $200\Delta_0+B_0=1$ and $50\Delta_0+B_0=1/3$, giving
\[
 \boxed{\Delta_0=\frac1{225},\quad B_0=\frac19,\quad V_0=\frac59}.
\]
At either node the old portfolio value equals the new portfolio cost, so the strategy is self-financing.

Now impose nonnegative holdings. At the low node, covering the payoff at price 100 requires $100\Delta+B\ge1$, hence the cost $50\Delta+B\ge1/2$. Holdings $\Delta=1/100,B=0$ attain this bound. At the high node, $100\Delta+B\ge1$ implies cost $200\Delta+B\ge1$, attained by cash 1. Initial wealth must therefore cover node values 1 and $1/2$. Every feasible initial portfolio satisfies
\[
 100\Delta_0+B_0
 =\tfrac13(200\Delta_0+B_0)+\tfrac23(50\Delta_0+B_0)
 \ge\frac23.
\]
Holdings $\Delta_0=1/300,B_0=1/3$ reach exactly those node values and permit the stated continuations. Thus
\[
 \boxed{V_0^{\rm super}=\frac23}.
\]
\insight Exact replication borrows at the low node. Once borrowing is prohibited, the cost must be recomputed by nodewise superhedging rather than taken from the unconstrained price.

\problem{R20}{Stopping with Recall of an Offer}{A limited observation budget}
Observe $X\sim U[0,1]$ for free. Either accept it, or pay $c>0$ to observe an independent $Y\sim U[0,1]$, after which you receive $\max(X,Y)$. Only one further observation is allowed. Find the optimal rule and initial net value. Evaluate $c=1/8$, and compare with the rule requiring acceptance of $Y$ after payment, with no recall of $X$.
\solution Given $X=x$, the expected gain from the additional observation rather than stopping is
\[
 \E(Y-x)^+-c=\int_x^1(y-x)\dd y-c=\frac{(1-x)^2}{2}-c.
\]
For $0<c<1/2$, continue exactly when $x<1-\sqrt{2c}$; either action is optimal at equality. For $c\ge1/2$, always stop. The optimal value is
\[
 V=\frac12+\int_0^{1-\sqrt{2c}}
 \left(\frac{(1-x)^2}{2}-c\right)\dd x
 =\frac23-c+\frac{(2c)^{3/2}}3
\]
for $c<1/2$, and $V=1/2$ otherwise. With $c=1/8$, the threshold is $1/2$ and $\boxed{V=7/12}$.

Without recall, the continuation value is the constant $\E Y-c=3/8$. Continue only for $x<3/8$. The net value is
\[
 \boxed{\E\max(X,3/8)=\frac12+\frac12\left(\frac38\right)^2
 =\frac{73}{128}}.
\]
Recall gives the higher value because the old offer remains a floor on the eventual prize.

\problem{R21}{Stopping while Learning a Win Rate}{Bayesian decisions}
An unknown parameter $P$ is drawn from $U[0,1]$. Conditional on $P$, successive coin tosses are independent with H probability $P$. Each H pays $+1$ and each T pays $-1$. You may stop at any time, including before the first toss, and may toss at most three times. Past rewards are retained. Find the maximum expected total reward and a complete optimal policy. Compare with a coin known to be fair.
\solution After $h$ H's and $t$ T's, the posterior density is proportional to $p^h(1-p)^t$, so the next H probability is $q_{h,t}=(h+1)/(h+t+2)$. Let $V_k(h,t)$ be the optimal additional reward with $k$ tosses left. Then $V_0=0$ and
\[
 \begin{aligned}
 C_k(h,t)&=q_{h,t}[1+V_{k-1}(h+1,t)]\\
 &\quad +(1-q_{h,t})[-1+V_{k-1}(h,t+1)],\\
 V_k(h,t)&=\max\{0,C_k(h,t)\}.
 \end{aligned}
\]
Stopping gives additional reward zero. The continuation term includes the next outcome and subsequent optimal decisions, so backward induction proves the recursion. With one toss left,
\[
 V_1(2,0)=\tfrac12,\qquad V_1(1,1)=V_1(0,2)=0.
\]
The preceding stages give
\[
 V_2(1,0)=\tfrac23,\quad V_2(0,1)=0,\qquad
 \boxed{V_3(0,0)=\tfrac13}.
\]
Toss initially. Stop after a first T and continue after a first H. After HH, continue once more. After HT, the final toss has conditional mean zero, so stopping and continuing tie. For a coin known to be fair, bounded optional stopping makes the expected reward zero under every rule using at most three tosses.
\insight Each toss has marginal H probability $1/2$, but observations change the predictive probability. A shared unknown parameter makes the outcomes dependent before conditioning on it.

\problem{R22}{Discounted Offers over an Unlimited Horizon}{A contraction argument}
Offers $X_1,X_2,\ldots$ are independent $U[0,1]$ variables. After each observation, accept it or discard it permanently. Acceptance at observation $N$ pays present value $\delta^{N-1}X_N$, where $0<\delta<1$. Never accepting pays zero. Find the optimal rule, initial value and mean observation count under a strict threshold rule.
\solution Let $V_k$ be the value with at most $k$ observations, and $V_0=0$. Backward induction gives
\[
 V_{k+1}=\E\max(X,\delta V_k).
\]
The map has Lipschitz constant at most $\delta<1$, so iteration converges to its unique fixed point $V$. Truncating any infinite-horizon rule at observation $k$ changes its payoff by at most $\delta^{k-1}$ in absolute value. Thus the fixed point is the infinite-horizon optimum.

Put $t=\delta V\in(0,1)$. Integration gives $V=(1+t^2)/2$, or $\delta t^2-2t+\delta=0$. Taking the root in $(0,1)$ yields
\[
 \boxed{t=\frac{1-\sqrt{1-\delta^2}}\delta},\qquad
 \boxed{V=\frac1{1+\sqrt{1-\delta^2}}}.
\]
Accept at the first $X_n>t$. The observation count is geometric of parameter $1-t$, hence $\boxed{\E N=1/(1-t)}$. Its actual discounted payoff is
\[
 \sum_{n=1}^\infty(\delta t)^{n-1}\E[X\ind_{\{X>t\}}]
 =\frac{1-t^2}{2(1-\delta t)}=V,
\]
verifying attainment.
\insight Discounting bounds the contribution of the distant future by a deterministic geometric tail, justifying finite-horizon approximation of the unlimited problem.

\problem{R23}{Choosing Which Die to Reroll}{Two-dimensional observations}
Roll two independent fair six-sided dice. After seeing both values, take their product, or pay $c\ge0$ to reroll exactly one chosen die once. The new value must be used; the payoff is the final product minus the fee. Find the complete optimal rule and the maximum expected net payoff for $c=1$.
\solution Sort the observed values as $x\le y$. Rerolling the smaller die has net expectation $(7/2)y-c$, whereas rerolling the larger gives $(7/2)x-c$. The former is at least as large. Thus
\[
 \boxed{\text{reroll the smaller die}\quad\Longleftrightarrow\quad
 y(7/2-x)>c},
\]
with either action optimal at equality; otherwise stop. Without intervention the product mean is $(7/2)^2=49/4$, so
\[
 V(c)=\frac{49}{4}+\frac1{36}\sum_{i=1}^6\sum_{j=1}^6
 \left(\tfrac72\max(i,j)-ij-c\right)^+.
\]
For $c=1$, continue exactly when the smaller value is at most 3. Grouping the 36 equally likely ordered outcomes by smaller value 1, 2 or 3 gives total gains $183/2,48,19/2$, respectively. Equal values count once and unequal pairs twice. Hence
\[
 \boxed{V(1)=\frac{49}{4}+\frac{149}{36}=\frac{295}{18}}.
\]
\insight Compare the two controlled actions first, then compare the better one with stopping. A low observed value alone does not justify a reroll without accounting for its fee.

\problem{R24}{Ordering Tests with Early Stopping}{Adjacent interchange}
Independent tests $i=1,\ldots,n$ pass with probabilities $p_i\in(0,1)$ and incur costs $c_i>0$. To determine whether all pass, choose their order, possibly adaptively, and stop at the first failure; if all pass, perform every test. Find an order minimizing expected cost. How does the rule change when the goal is to determine whether at least one passes? Evaluate $p=(9/10,1/2,1/10)$ and $c=(1,2,5)$.
\solution For the all-pass decision, every nonterminal history consists entirely of passes. Thus each deterministic adaptive policy is equivalent to a fixed permutation; randomized policies are mixtures of permutations. The expected cost of permutation $\pi$ is
\[
 \sum_{k=1}^n c_{\pi_k}\prod_{j<k}p_{\pi_j}.
\]
Adjacent tests $i,j$ contribute $c_i+p_ic_j$ or $c_j+p_jc_i$. The probability of continuing past both is $p_ip_j$ in either order, so later costs are unchanged. Putting the cheaper order first is equivalent to
\[
 c_i(1-p_j)\le c_j(1-p_i)
 \quad\Longleftrightarrow\quad
 \frac{c_i}{1-p_i}\le\frac{c_j}{1-p_j}.
\]
Successively removing inversions proves that increasing $c_i/(1-p_i)$ is globally optimal.

For the at-least-one-pass decision, stop at the first pass; every surviving history consists of failures. Replacing the continuation probabilities by $1-p_i$ in the same interchange proof gives increasing $c_i/p_i$. For the stated data,
\[
 \boxed{\text{all pass: }(2,3,1),\quad C=\frac{91}{20}},
\]
\[
 \boxed{\text{at least one: }(1,2,3),\quad C=\frac{29}{20}}.
\]
\insight The ordering index depends on the logical stopping objective. The same tests can require different optimal orders for different decisions.

\par\Needspace{8\baselineskip}
\subsection*{Section Exercises}
\exercise{OR20}{Optimal Stopping with Red and Black Cards}
A deck contains 26 red cards and 26 black cards and is uniformly shuffled. Cards are drawn one at a time without replacement. The score starts at 0, increases by 1 for a red card, and decreases by 1 for a black card. You may stop before drawing or after any draw and receive the current cumulative score. Find an optimal stopping rule and the maximum expected score.

\exercise{OR22}{Rerolling a Die at a Cost}
A fair six-sided die has faces $1,2,\ldots,6$. The first roll is free. After observing a face, you may accept it and stop, or pay 1 to roll again; every subsequent reroll also costs 1. Rolls are independent, and previous faces cannot be recalled. Admissible strategies must stop after finitely many rolls almost surely. The final payoff is the last face minus all reroll fees. Find an optimal strategy and the maximum expected payoff.


\section{Brownian Motion}
\noindent\textbf{Key Ideas and Methods.}\enspace
Independent increments and Gaussian structure reduce many Brownian calculations to covariance integrals. Deterministic square-integrable kernels give Gaussian It\^o integrals; exponentiation or evaluation at an independent random time generally does not preserve normality. Use orthogonal projection for jointly Gaussian observations and reflection for path maxima. Conditioning on an endpoint, a probability-zero event, requires densities or a continuous version of the conditional law.

An independent exponential failure time becomes a discount factor; a generator boundary problem can then be verified with It\^o's formula. Interchanging integrals, expectations or parameter derivatives requires nonnegativity, integrability or a suitable domination argument. Poisson splitting requires independent marking. In the renewal model, uniform observation over a long time interval asymptotically selects intervals in proportion to their lengths.
\par\smallskip



\problem{R25}{Predicting an Exponential Brownian Integral}{A non-Gaussian projection}
Let $W$ be standard Brownian motion starting at 0, let $T>0$, and define
\[
 I=\int_0^T e^{W_t-t/2}\dd t.
\]
Find its mean and variance, and the coefficients and error of the best predictor $\alpha+\beta W_T$ under squared loss. Then find the conditional mean given $W_T=b$ and determine whether it equals that affine predictor.
\solution Write $M_t=e^{W_t-t/2}$. The normal moment generating function gives $\E M_t=1$ and, for $s\le t$, $\E M_sM_t=e^s$. Tonelli's theorem for the nonnegative integrands yields
\[
 \E I=T,\qquad
 \E I^2=2\int_0^T\int_0^t e^s\dd s\dd t
 =2(e^T-1-T).
\]
Finiteness also supplies the square integrability needed for subsequent covariance interchanges. Hence
\[
 \boxed{\Var I=2(e^T-1-T)-T^2}.
\]
Independent increments and differentiation of the normal moment generating function give $\E(M_tW_T)=\E(M_tW_t)=t$, so $\Cov(I,W_T)=T^2/2$. The normal equations for squared loss give
\[
 \boxed{\alpha=T,\qquad\beta=T/2}.
\]
\[
 \boxed{\operatorname{MSE}_{\rm affine}=2(e^T-1-T)-T^2-\frac{T^3}{4}}.
\]
Given $W_T=b$, $W_t$ is normal with mean $tb/T$ and variance $t-t^2/T$. This follows from the Gaussian orthogonal decomposition of $W_t-(t/T)W_T$ and $W_T$. Using the normal exponential moment and conditional Tonelli again,
\[
 \boxed{\E(I\mid W_T=b)=\int_0^T
 \exp\left(\frac{tb}{T}-\frac{t^2}{2T}\right)\dd t}.
\]
This function has strictly positive second derivative in $b$, so it is not affine. Allowing nonlinear prediction strictly reduces mean squared error.
\insight Exponentiation destroys the joint Gaussian structure. Covariance still determines the best linear coefficient, but it does not force the conditional mean to be linear.

\problem{R26}{The Maximum Given a Brownian Endpoint}{Reflection principle}
Let $W$ be standard Brownian motion and $M_T=\max_{0\le t\le T}W_t$, where $T>0$. Find the conditional law of $M_T$ given $W_T=b$. In particular, find its mean when $b=0$.
\solution For $a>\max(0,b)$, reflect the path after its first visit to $a$ about that level. Paths that reach $a$ and finish near $b$ correspond to paths ending near $2a-b$. Writing $\varphi_T$ for the $N(0,T)$ density, the conditional tail is
\[
 \boxed{\Prob(M_T\ge a\mid W_T=b)
 =\frac{\varphi_T(2a-b)}{\varphi_T(b)}
 =\exp\left(-\frac{2a(a-b)}T\right)}.
\]
This density ratio defines a continuous version of the regular conditional law; it is not a ratio of probabilities of zero-probability events. For $a\le\max(0,b)$, the tail probability is 1.

When $b=0$, the maximum is nonnegative and its distribution function is $1-e^{-2a^2/T}$ for $a\ge0$. The tail integral gives
\[
 \boxed{\E(M_T\mid W_T=0)=\int_0^\infty e^{-2a^2/T}\dd a
 =\sqrt{\frac{\pi T}{8}}}.
\]
\insight An endpoint constraint does not constrain the whole path. Even a path ending at zero may have crossed a positive barrier; reflection computes that crossing probability.

\problem{R27}{Conditioning on Both Endpoint and Area}{Gaussian conditioning}
Let $W$ be standard Brownian motion, $T>0$, and $A=\int_0^T W_t\dd t$. Conditional on $W_T=b$ and $A=a$, find the mean and covariance functions of the entire process. Explain why two path values can have negative conditional covariance.
\solution First remove the endpoint direction by setting $B_t=W_t-(t/T)W_T$. Joint Gaussianity makes $B$ independent of $W_T$, with covariance
\[
 K(s,t)=\Cov(B_s,B_t)=\min(s,t)-\frac{st}{T}.
\]
Next put $J=\int_0^T B_t\dd t=A-TW_T/2$. The bounded continuous covariance ensures square integrability of the integral. Integrating the covariance gives
\[
 \Cov(B_t,J)=\frac{t(T-t)}2,\qquad \Var J=\frac{T^3}{12}.
\]
Project $B_t$ onto $J$. The residual is independent of both $J,W_T$, jointly for every finite selection of times. Thus the conditional process remains Gaussian, with mean
\[
 \boxed{m(t)=\frac{tb}{T}
 +\frac{6t(T-t)}{T^3}\left(a-\frac{Tb}{2}\right)},
\]
and covariance
\[
 \boxed{K_*(s,t)=\min(s,t)-\frac{st}{T}
 -\frac{3s(T-s)t(T-t)}{T^3}}.
\]
The two observations have a nonsingular covariance matrix because $\Var J>0$, so these formulas define a continuous conditional version for every real $a,b$. Taking $s=T/4,t=3T/4$ gives $K_*(s,t)=-11T/256<0$. Fixing total area allows an unusually high value at one time to be compensated by a lower value elsewhere, unlike the original positive Brownian covariances.

\problem{R28}{Hitting a Boundary before Independent Failure}{A discounted boundary problem}
Let $X_t=x+\sigma W_t$, where $0<x<L$ and $\sigma>0$, and let $\tau$ be its first hit of 0 or $L$. An independent exponential failure clock $\zeta$ has rate $\lambda>0$. Success means reaching $L$ before both 0 and failure. Find the success probability and the mean duration conditional on success, including their limits as $\lambda\downarrow0$.
\solution Independence of the failure clock gives
\[
 u_\lambda(x)=\E_x[e^{-\lambda\tau}\ind_{\{X_\tau=L\}}].
\]
Consider the boundary problem
\[
 \frac{\sigma^2}{2}u''-\lambda u=0,\qquad u(0)=0,\quad u(L)=1.
\]
Put $k=\sqrt{2\lambda}/\sigma$. The general solution is spanned by $\sinh(kx),\cosh(kx)$, and the endpoints determine
\[
 \boxed{u_\lambda(x)=\frac{\sinh(kx)}{\sinh(kL)}}.
\]
To verify the interpretation, apply It\^o's formula to $e^{-\lambda(t\wedge\tau)}u(X_{t\wedge\tau})$; its drift is zero. Boundedness of $u,u'$ on the interval makes the stochastic integral mean zero at each finite time. A Brownian increment large enough to cross the interval has a uniform positive probability over a fixed time block, so $\tau<\infty$ almost surely. Dominated convergence gives the discounted terminal-payoff identity.

For $\lambda>0$, $\tau e^{-\lambda\tau}$ is bounded over nonnegative times, uniformly for $\lambda$ in a sufficiently small neighborhood. Differentiation under the expectation is therefore justified, giving
\[
 \E_x(\tau\mid\text{success})=-\partial_\lambda\log u_\lambda(x)
 =\boxed{\frac{L\coth(kL)-x\coth(kx)}{\sigma^2k}}.
\]
Using $\sinh z\sim z$ and $\coth z=1/z+z/3+O(z^3)$,
\[
 \boxed{u_\lambda(x)\to\frac xL},\qquad
 \boxed{\E_x(\tau\mid\text{success})\to\frac{L^2-x^2}{3\sigma^2}}.
\]
\insight Independent failure weights a path of duration $\tau$ by $e^{-\lambda\tau}$, favoring faster successes. The negative logarithmic derivative of success probability yields the conditional duration.

\problem{CR01}{Poisson Splitting and a First Arrival}{Classical example}
Events arrive as a Poisson process of rate $\lambda>0$. Independently label each arrival A with probability $p$ and B with probability $q=1-p$, where $0<p<1$. Write $N_A(t),N_B(t)$ for the two counts. Prove that they are independent Poisson processes, and find the A count conditional on a total of $n$ arrivals. Then let $\tau$ be the first A arrival and $K=N_B(\tau)$. Find the distribution of $K$ and the conditional distribution of $\tau$ given $K=k$.
\solution Conditioning first on the total count, the joint probability generating function at a fixed time is
\[
 \E[z^{N_A(t)}w^{N_B(t)}]
 =\exp\{\lambda t(pz+qw-1)\}
 =e^{\lambda pt(z-1)}e^{\lambda qt(w-1)}.
\]
It factors into generating functions for Poisson counts with rates $\lambda p$ and $\lambda q$. Original increments over disjoint intervals are independent, as are all labels. The same factorization therefore holds jointly over any finite collection of disjoint intervals, proving independence of the two processes. Given a total of $n$ arrivals, their independent labels directly give
\[
 \boxed{N_A(t)\mid N_A(t)+N_B(t)=n\sim\operatorname{Bin}(n,p)}.
\]
Put $\alpha=\lambda p,\beta=\lambda q$. The first A arrival has distribution $\operatorname{Exp}(\alpha)$ and is independent of the entire B process. Hence
\[
 \Prob(K=k,\tau\in\dd t)
 =\alpha e^{-(\alpha+\beta)t}\frac{(\beta t)^k}{k!}\dd t,
 \qquad k=0,1,\ldots.
\]
Using $\int_0^\infty t^ke^{-\lambda t}\dd t=k!/\lambda^{k+1}$ gives
\[
 \boxed{\Prob(K=k)=pq^k}.
\]
Dividing the joint density by this probability yields
\[
 \boxed{f_{\tau\mid K=k}(t)=\frac{\lambda^{k+1}}{k!}t^ke^{-\lambda t},\quad t>0},
 \qquad \boxed{\E(\tau\mid K=k)=\frac{k+1}{\lambda}}.
\]
This is the sum of $k+1$ independent exponential interarrival times of rate $\lambda$: exactly $k$ B arrivals precede the first A.
\insight Independence of counts at fixed times does not imply independence after stopping at the first arrival of one type. A joint density makes the resulting conditioning explicit.

\problem{CR02}{Residual Life at a Random Time}{Classical example}
A renewal process starts at time zero. Independent intervals have lengths 1 and 3, each with probability $1/2$. Independently choose $U_t$ uniformly in $[0,t]$. Let $L_t$ be the length of the complete renewal interval containing $U_t$, $A_t$ the time since its left endpoint and $R_t$ the time to its right endpoint. Find their joint limiting distribution as $t\to\infty$. Under that limit, compute $\E R$, $\Prob(R>2)$ and $\Cov(A,R)$.
\solution A uniformly chosen time selects intervals in proportion to their duration. More precisely, among the first $n$ complete intervals, the fraction of total time spent in intervals of length $\ell$ converges by the strong law to
\[
 \frac{\ell\Prob(L=\ell)}{\E L},\qquad \E L=2.
\]
The unfinished interval at the right endpoint has length at most 3, so its contribution to the time fraction vanishes. The limiting selected length $L_*$ therefore satisfies
\[
 \boxed{\Prob(L_*=1)=\frac14,\qquad\Prob(L_*=3)=\frac34}.
\]
Within an interval of length $\ell$, a uniformly selected time has a uniform location. Conditional on $L_*=\ell$,
\[
 \boxed{A\sim U[0,\ell],\qquad R=\ell-A}.
\]
This mixture specifies the joint distribution. All variables are bounded, so the distributional limit also gives the required moment limits. Since $\E L^2=5$ and $\E L^3=14$,
\[
 \boxed{\E R=\E A=\frac{\E L^2}{2\E L}=\frac54},\qquad
 \boxed{\Prob(R>2)=\frac{\E(L-2)^+}{\E L}=\frac14}.
\]
The integral of the product of age and residual life over an interval of length $\ell$ is $\int_0^\ell u(\ell-u)\dd u=\ell^3/6$. Consequently,
\[
 \E(AR)=\frac{\E L^3}{6\E L}=\frac76,
 \qquad \boxed{\Cov(A,R)=\frac76-\left(\frac54\right)^2=-\frac{19}{48}}.
\]
\insight Sampling an interval uniformly and sampling the interval containing a uniform time are different experiments. The latter first weights by length, then selects a uniform location within that interval.

\problem{CR03}{A Single-Server Busy Period}{Classical example}
Customers arrive as a Poisson process of rate $\lambda$. Independent service times are exponential with rate $\mu$, independently of arrivals, and $0<\lambda<\mu$. One server works in first-come, first-served order. Starting with one customer, let $B$ be the first time the system is empty. Find the Laplace transform, mean and variance of $B$, and explain what changes in the mean when $\lambda=\mu$.
\solution Regard arrivals during a customer's service as that customer's offspring. Conditional on the service time $S$, the offspring count is $\operatorname{Pois}(\lambda S)$, with mean $\rho=\lambda/\mu<1$. The expected size of generation $j$ is $\rho^j$, so the expected total population is $1/(1-\rho)$. Its finiteness implies that the total population is finite almost surely. Existence of a customer does not depend on that customer's own future service time. Tonelli's theorem gives
\[
 \E B=\frac1\mu\sum_{j=0}^\infty\rho^j=\frac1{\mu-\lambda}.
\]
Let $b(s)=\E e^{-sB}$ for $s\ge0$. Arrivals during the first service generate independent copies of the descendant service total. The busy period is the sum of all service times in the tree, regardless of processing order. Conditioning on $S$ therefore gives
\[
 b(s)=\E\left[e^{-sS}b(s)^{N(S)}\right]
 =\frac{\mu}{\mu+s+\lambda[1-b(s)]}.
\]
Thus $\lambda b^2-(\lambda+\mu+s)b+\mu=0$. Choose the root in $(0,1)$ for $s>0$ and extend continuously to zero:
\[
 \boxed{b(s)=\frac{\lambda+\mu+s-\sqrt{(\lambda+\mu+s)^2-4\lambda\mu}}{2\lambda}}.
\]
Differentiation for $s>0$ is valid because $B^je^{-sB}$ is bounded. Differentiating the quadratic equation twice and writing $D_s=2\lambda b(s)-\lambda-\mu-s$ yields
\[
 D_s b'(s)=b(s),\qquad
 D_s b''(s)=2b'(s)-2\lambda[b'(s)]^2.
\]
As $s\downarrow0$, $b(s)\to1$ and $D_s\to\lambda-\mu$. Monotone convergence for the nonnegative moment integrands identifies the finite derivative limits:
\[
 \boxed{\E B=\frac1{\mu-\lambda}},\qquad
 \E B^2=\frac{2\mu}{(\mu-\lambda)^3},
\]
\[
 \boxed{\Var B=\frac{\lambda+\mu}{(\mu-\lambda)^3}}.
\]
At criticality, the offspring generating function is $f(z)=\mu/[\mu+\lambda(1-z)]$. Iteration of $f$ from zero gives extinction probabilities within finitely many generations; the limit is the smallest fixed point in $[0,1]$. When $\lambda=\mu$, that fixed point is 1. The busy period still ends almost surely, but every generation has expected size 1. The nonnegative series for expected total service therefore diverges, giving $\E B=\infty$.
\insight A busy period includes every generation of arrivals during service. Counting only the initial service and its immediate arrivals omits later generations.


\par\Needspace{8\baselineskip}
\subsection*{Section Exercises}
\exercise{OR25}{Correlation of Brownian Integrals}
Let $(W_t)_{t\ge0}$ be standard Brownian motion and let $T>0$ be deterministic. Define
\[
 X=\int_0^T W_t\,dt,\qquad
 Y=\int_0^T t\,dW_t,
\]
where the second integral is an It\^o integral. Find the correlation of $X$ and $Y$.

\exercise{OR26}{Where Planar Brownian Motion First Hits an Axis}
Standard planar Brownian motion starts at $(1,1)$, so its position is $(1+B_t^{(1)},1+B_t^{(2)})$, where the two coordinates are independent standard one-dimensional Brownian motions. Stop when the path first hits the $y$-axis. Find the probability that the stopping point lies on its positive half-axis.


\section{Exercises}
\noindent\textbf{Key Ideas and Methods.}\enspace
Start by fixing the counting conventions: initial visits, the terminal step, charges and self-loops at an empty state. Then identify the needed prefix, direction, posterior parameter or conditional covariance, rather than selecting a formula from the problem's wording.

Check candidate values against boundaries, normalization and feasible actions. In sums stopped at a random time, determine which continuation events depend only on past information. Conditional durations and distributions require the appropriate reweighting; unconditional formulas need not remain valid.
\par\smallskip



\exercise{RS1}{Waiting for a Self-Overlapping Word}
Toss a fair coin independently until substring HTHT first appears. Starting with an empty record, use prefix states to find the mean number of tosses. Explain which state survives when H follows HTH.
\exercise{RS2}{A Clearing Queue's Tail}
Each queue operation independently adds one item with probability $1/4$, serves one with probability $1/2$ (doing nothing if empty), or clears the queue with probability $1/4$. Find the stationary probability of length at least 2 and the mean positive-time return to an empty queue starting empty.
\exercise{RS3}{Subsampling and Identifiability}
A symmetric two-state Markov chain switches with unknown one-step probability $\alpha\in(0,1)$. Its two-step switch probability is known exactly to be $3/8$. Find all possible values of $\alpha$. Can observations at even times alone distinguish them?
\exercise{RS4}{Variance of Recursive Work}
An A task produces two B tasks with probability $1/2$ and none otherwise. A B task produces one A with probability $1/3$ and none otherwise. Choices are independent. Starting from one A, find the mean, second moment and variance of the total task count, and justify finiteness.

\exercise{RH1}{Restart Counts and Costs}
An integer process starts at 1 and stops at 0. Each step resets to 1 with probability $1/2$; otherwise it moves one unit left or right with equal probability. A reset costs 3 and a move costs 1. Find the mean duration, reset count and total cost.
\exercise{RH2}{A Three-Unit Drawdown}
A fair integer walk starts at 0 and stops when its drawdown from the running maximum first reaches 3. Find the mean and variance of the time, and the probability that its maximum before stopping is at least 4.
\exercise{RH3}{Information in the Initial Direction}
A persistent walk on $\{0,\ldots,6\}$ stops at either endpoint. Each step keeps the preceding direction with probability $3/4$ and reverses otherwise. Starting from 3, find the probability of reaching 6 first for a rightward and a leftward preceding direction. What if that unobserved initial direction is equally likely to be either?
\exercise{RH4}{Costs until a Repeat}
Independent draws have three equally likely types with costs 1, 2 and 4. Stop at the first repeated type, charging for the final draw as well. Find the mean total cost and the mean number of distinct types seen by stopping.

\exercise{RC1}{A Put Payoff without Short Selling}
A stock costs 1 and is worth 2 or $1/2$ one period later. Cash earns zero interest and fractional holdings are allowed. The payoff is $(1-S_1)^+$. Find the replication price and holdings when borrowing and short selling are allowed. If stock holdings must be nonnegative, find the smallest initial wealth covering the payoff on every path.
\exercise{RC2}{Learning with Two Opportunities}
A coin's unknown H probability has prior $U[0,1]$; conditional on it, tosses are independent. H pays $+1$ and T pays $-1$. Past earnings are kept. You may toss at most twice and may stop at any time. Find the optimal rule and initial value.
\exercise{RC3}{Replacing a Correlated Offer}
Receive $X\sim U[0,1]$ for free. Accept it, or pay $1/8$ to replace it by $1-X$. Only one replacement is allowed, with no recall. The relationship between the offers is known. Find the optimal rule and expected net payoff.
\exercise{RC4}{Two Testing Objectives}
Three independent tests have pass probabilities $(1/2,3/4,1/4)$ and costs $(1,2,1)$. Find the optimal order and expected cost for deciding whether all pass, and for deciding whether at least one passes. Stop as soon as the required conclusion is known.

\exercise{RG1}{Quadratic Prediction of an Exponential Integral}
For standard Brownian motion, let $I=\int_0^1 e^{W_t-t/2}\dd t$. Find the predictor minimizing mean squared error among $\alpha+\beta W_1+\gamma(W_1^2-1)$, and its error.
\exercise{RG2}{Negative Covariance under Two Constraints}
For standard Brownian motion, put $A=\int_0^2W_t\dd t$. Condition on $W_2=0$ and $A=1$. Find the conditional law of $W_1$ and the conditional covariance of $W_{1/2},W_{3/2}$.
\exercise{RG3}{A Bridge's Maximum}
Condition standard Brownian motion on $W_2=0$. Find the conditional probability that $\max_{0\le t\le2}W_t\ge1$ and the conditional mean of this maximum.
\exercise{RG4}{A Successful Hit before Failure}
Standard Brownian motion starts at 1 and stops at 0 or 2. An independent failure time is exponential with rate $1/2$. Find the probability of reaching 2 before failure and before 0, and the mean duration conditional on success. Compare with the absence of a failure clock.
